\documentclass{article}
\usepackage[left=2cm, right=2cm, top=2cm, bottom=2cm]{geometry}
\usepackage{graphicx}
\usepackage{amsmath}
\usepackage{amssymb}
\usepackage{amsthm}
\usepackage{fancyhdr}
\usepackage{verbatim}
\usepackage{listings}
\usepackage{xcolor}
\usepackage{pdfpages}
\usepackage{cancel}
\usepackage{physics}
\usepackage{esint}

\lstset{
    language=C++,
    basicstyle=\ttfamily\footnotesize,
    keywordstyle=\color{blue},
    commentstyle=\color{green},
    stringstyle=\color{red},
    numbers=left,
    numberstyle=\tiny\color{gray},
    frame=single,
    breaklines=true,
    captionpos=b,
}


\title{PHYS 487: Lecture \# 18}
\author{Cliff Sun}

\newtheorem{theorem}{Theorem}[section]
\newtheorem{lemma}[theorem]{Lemma}
\newtheorem{definition}[theorem]{Definition}
\newtheorem{conjecture}[theorem]{Conjecture}
\newtheorem{proposition}[theorem]{Proposition}
\newtheorem{corollary}[theorem]{Corollary}
\newtheorem{one minute paper}[theorem]{One Minute Paper}

\pagestyle{fancy}
\lhead{\textbf{\thepage}\ \ \nouppercase{\rightmark}}
\chead{PHYS 487: Lecture \# 18}
\rhead{Cliff Sun}

\begin{document}

\maketitle

\section*{Lecture Span}
\begin{enumerate}
    \item Adiabatic Theorem
\end{enumerate}

\section*{Recap}

Consider a sinusoidal perturbation, then also consider an initial state $\ket{i}$ and final state $\ket{f}$. Where $\ket{f}$ is a continuum of states. Then the probability of leaving $\ket{i}$ is 

\begin{equation}
    \int dE_n \rho(E_n)\frac{|\langle i | V | n \rangle |^2}{\hbar}\frac{\sin^2(\Delta \omega t/2)}{\Delta \omega^2}
\end{equation}

Then, the rate of leaving $\ket{i}$ is 

\begin{equation}
    R = \frac{2\pi}{\hbar}\frac{|V_{if}|^2}{2}\rho(E_f)
\end{equation}

Consider a chunk of material, now shoot some light into it along the $z$ direction. Which direction does the electron most likely come out? Let 

$$\psi_0(r) = \frac{1}{\sqrt{\pi a^3}}e^{-r/a}$$

Firstly, find $\rho$. Then also find $V_{if}$. 

\section*{Adiabatic Theorem \& Approximation}

Adiabatic in stat mech is that no heat or mass exchange with environment. But quantum mechanically, if we are in a stationary state, then we will not leave it. 

For instance, choose $$\ket{1} = \sqrt{\frac{2}{L}}\sin(\pi x/L), \quad E_1 = \frac{\pi^2 \hbar^2 }{2m L^2}$$

If we extend the ground of the well to $2L$, then if the adiabatic theorem is true, then we will still be in the ground state. However, the final energy will be reduced by $1/4$. 
This is different. The rest of the energy is put into making the well larger. 

If we do this faster, then we can expand out the initial wavefunction in the basis states of the new well. But this requires a gap between the ground state and the excited states as a function of time. 

\begin{proof}
    Probability of leaving the state when changing the Hamiltonian. At all times 
    \begin{equation}
        H(t)\ket{n(t)} = E_n(t)\ket{n(t)}
    \end{equation}
    Time independent schrodinger wave equation for all $t$. Any state $\ket{\psi(t)} = \sum_n c_n(t)\ket{n(t)}$.
    \begin{equation}
        i\hbar \ket{\dot{\psi}(t)} = H(t)\ket{\psi(t)}
    \end{equation}
    Then, 
    \begin{equation}
        i\hbar \partial_t\left( \sum_n c_n(t)\ket{n(t)} \right) = H(t)\sum_n c_n(t)\ket{n(t)}
    \end{equation}
    \begin{equation}
        i\hbar \left( \sum_n \left( \dot{c}_n\ket{n} + c_n\ket{\dot{n}} \right) \right) = H(t)\sum_n c_n \ket{\psi}
    \end{equation}
    Then, multiply by $\bra{m(t)}$ and obtain 
    \begin{equation}
        i\hbar \left[ \dot{c}_m(t) + \sum_n \langle m | c_n | \dot{n} \rangle  \right]= \sum_n c_n(t) \langle m(t) | H(t) | n(t) \rangle
    \end{equation}
    \begin{equation}
        i\hbar\left( \dot{c}_m(t) + \sum_n  \langle m(t) | \dot{n}(t) \rangle \right) = c_m(t)E_m(t)
    \end{equation}
    Looking at $H(t)\ket{n(t)} = E_n(t)\ket{n(t)}$. Then take the time derivative:
    \begin{equation}
        \dot{H}n + H\dot{n} = \dot{E}_nn + E_n\dot{n}
    \end{equation}
    Then, 
    \begin{equation}
        m \dot{H} n + E_m m \dot{n} = E_n(t)\delta_{mn} + E_n(t) m \dot{n}
    \end{equation}
    For $n \neq m$, then 
    \begin{equation}
        \langle m | \dot{n} \rangle = \frac{\langle m | \dot{H} | n \rangle}{E_n - E_m}
    \end{equation}
    For $m = n$, then 
    \begin{equation}
        \langle m(t) | \dot{H} | m(t) \rangle = \dot{E}(t)
    \end{equation}
    Insert this back into our original equation: 
    \begin{equation}
        \dot{c}_m(t) + \left( \frac{i}{\hbar}E_m(t) + \langle m(t) | \dot{m}(t) \rangle \right)c_m(t) = \sum_{n \neq m} \frac{\langle m | \dot{H} | n \rangle}{E_m - E_n}c_n(t)
    \end{equation}
    Adiabatic approximation: assume that 
    \begin{equation}
        \frac{\langle m | \dot{H} | n \rangle}{E_m - E_n} \ll \frac{|E_n(t) - E_m(t)|}{\hbar}, \quad \forall t
    \end{equation}
    This is saying that the frequency generated by evolving the Hamiltonian does not excite all states. Then 
    \begin{equation}
        \dot{c}_m(t) = \left( -\frac{i}{\hbar}E_m(t) + \langle m(t)| \dot{m}(t) \rangle \right)c_m(t)
    \end{equation} 
    After integration, 
    \begin{equation}
        c_m(t) = c_m(0)e^{i\Theta_m(t)}e^{i\gamma_m (t)}
    \end{equation}
    Where, (dynamic phase)
    \begin{equation}
        \Theta_m(t) = \int\left( -\frac{i}{\hbar}E_m(t)\right)
    \end{equation}
    And (geometric phase)
    \begin{equation}
        \gamma_m(t) = i\int\left( \langle m(t)| \dot{m}(t) \rangle \right)
    \end{equation}
    Here, $\gamma_m(t)$ is always real. Therefore, 
    \begin{equation}
        c_m(t)^2 = \text{const.}
    \end{equation}
    We stay in the initial eigenstate. 
\end{proof}

\section*{Two Level System}  

\begin{equation}
    H = \frac{\hbar(\omega - \omega_0)}{2}\left( \ket{0}\bra{0} - \ket{1}\bra{1} \right) + \frac{\hbar \Omega}{2}\left( \ket{0}\bra{1} + \ket{1}\bra{0} \right)
\end{equation}

Diagonalize this: 

\begin{equation}
    E_{\pm} = \frac{\hbar}{2}\sqrt{\Delta \omega^2 + \Omega^2}
\end{equation}

With eigenstates 

\begin{equation}
    \ket{\psi_{+}} = \cos\theta/2 \ket{0} + \sin\theta/2\ket{1}
\end{equation}
\begin{equation}
    \ket{\psi_{-}} = -\sin\theta/2 \ket{0} + \cos\theta/2\ket{1}
\end{equation}

Look at the energy as a function of $\Delta \omega$. 

\end{document}